\documentclass[../../main.tex]{subfiles}

\graphicspath{{../../image/}{../../figure/}}

\begin{document}

\section{平衡律方程组}\label{section:ch01-平衡律方程组}

我们考虑 $n$ 条不同的平衡律同时在 $\mathscr{X}$ 上作用的情形，其生成测度由生成密度场诱导；并把它们的场方程\cref{align:ch01-divA-P} 合并为方程组
\begin{align}
\operatorname{div} A(X)=P(X),\label{align:ch01-system-divA}
\end{align}
其中现在 $A$ 是 $n\times k$ 矩阵场，$P$ 是 $n$ 列向量场。散度算子作用于 $A$ 的行向量，结果为一个 $n$ 列向量场。

我们假定介质的状态由一个\emph{状态向量场} $U$ 描述，取值于 $\mathbb{R}^{n}$ 的一个开子集 $\mathcal{O}$，它通过\emph{本构方程}确定点 $X\in\mathscr{X}$ 处的通量密度场 $A$ 与生成密度场 $P$：
\begin{align}
A(X)=G(U(X),X),\qquad P(X)=\Pi(U(X),X),\label{align:ch01-constitutive}
\end{align}
其中 $G$ 与 $\Pi$ 是定义在 $\mathcal{O}\times\mathscr{X}$ 上的已知光滑函数，取值分别在 $\mathbb{M}^{n\times k}$ 与 $\mathbb{R}^{n}$ 中。

将\cref{align:ch01-system-divA} 与\cref{align:ch01-constitutive} 联立得到
\begin{align}
\operatorname{div} G(U(X),X)=\Pi(U(X),X),\label{align:ch01-closed-system}
\end{align}
即一个（形式上）封闭的拟线性一阶偏微分方程组，状态向量场须由它确定。形如\cref{align:ch01-closed-system} 的方程，当 $n\geqslant 2$ 时此后称为\emph{平衡律方程组}，当 $n=1$ 时称为\emph{标量平衡律}。在无生成即 $\Pi\equiv 0$ 的特殊情形，当 $n\geqslant 2$ 时称为\emph{守恒律方程组}，当 $n=1$ 时称为\emph{标量守恒律}。这一术语并非完全标准：某些作者宁愿称“带源守恒律方程组”。当 $G$ 与 $\Pi$ 不显式依赖于 $X$ 时，平衡律方程组称为\emph{齐次}的。

注意，当坐标在某固定点 $\bar{X}\in\mathscr{X}$ 附近拉伸，即 $X=\bar{X}+\varepsilon Y$ 时，随着 $\varepsilon\downarrow 0$，平衡律方程组\cref{align:ch01-closed-system} 化为关于变量 $Y$ 的齐次守恒律方程组。这就是为什么一般平衡律方程组解的局部性质可不失一般性地在齐次守恒律方程组这一更简单的框架中研究。

在 $\mathscr{X}$ 上几乎处处满足\cref{align:ch01-closed-system} 的 Lipschitz 连续场 $U$ 称为\emph{古典解}。一个可测场 $U$ 若在分布意义下满足\cref{align:ch01-closed-system}，即 $G(U(X),X)$ 与 $\Pi(U(X),X)$ 局部可积，且对任一检验函数 $\varphi\in C_{0}^{\infty}(\mathscr{X})$ 有
\begin{align}
\int_{\mathscr{X}}\bigl[G(U(X),X)\operatorname{grad}\varphi(X)+\varphi(X)\Pi(U(X),X)\bigr]\,\mathrm{d}X=0,\label{align:ch01-weak-system}
\end{align}
则称为\emph{弱解}。任何 Lipschitz 连续的弱解必为古典解。

\textbf{记号}：对 $\alpha=1,\cdots,k$，$G_{\alpha}(U,X)$ 表示矩阵 $G(U,X)$ 的第 $\alpha$ 列向量。

\medskip

\textbf{记号}：此后 $\mathsf{D}$ 表示关于变量 $U$ 的微分。与矩阵记号连用时，$\mathsf{D}$ 视为行算子：$\mathsf{D}=[\partial/\partial U^{1},\cdots,\partial/\partial U^{n}]$。

\end{document}
